% PLANNED-VALUE: every JAM score in this table is a layout target.
\begin{table}[t]
\centering\small
\begin{tabular}{lr}
\toprule
Planned capacity or annotation setting & LIBERO-Plus success \\
\midrule
Agent width 768; full action labels & \target{85.0} \\
Agent width 1024; full action labels & \target{90.2} \\
Agent width 1280; full action labels & \target{91.0} \\
Action-label density 25 percent & \target{82.0} \\
Action-label density 50 percent & \target{86.0} \\
Action-label density 100 percent & \target{90.2} \\
\bottomrule
\end{tabular}
\caption{\textbf{Which scale axis matters?} Layout targets in percent. Agent-width comparisons keep the video backbone, head geometry, depth, data, and update budget fixed and report the resulting compute separately. Annotation-density comparisons keep all video windows and robot exposure fixed while masking action labels at the episode level. The full-label setting is the shared reference for the density sweep. These settings require dedicated source configurations before execution. \targetnote}
\label{tab:scalecontrols}
\end{table}
